% Modul Belajar 9, dihasilkan oleh tools/build.py dari tools/konten/p09.py
\documentclass[a4paper]{article}
\usepackage{modulITERA}
\begin{document}
\Sampul{9}{Array Satu Dimensi}{Menyimpan banyak nilai sejenis, mengaksesnya dengan indeks, dan memprosesnya dengan perulangan}{CPMK0607 (menerapkan) dan CPMK0608 (menjelaskan) pada konsep dasar algoritma dan sistem komputer}{sekitar 3 jam belajar mandiri, ditambah 150 menit di kelas}{\texttt{02 Hands-on/P9 - Array 1D - Hands-on/}}
\Bagian{Modul 9}{Peta belajar}
Ikuti urutan ini dari atas ke bawah. Setiap bagian memakai apa yang sudah dibahas di bagian sebelumnya.
\par\vspace{4pt}\noindent{\fontsize{9.5pt}{13pt}\selectfont
\begin{tabular}{@{}>{\raggedright\arraybackslash}p{0.140\textwidth}>{\raggedright\arraybackslash}p{0.840\textwidth}@{}}
\kepala{Urutan} & \kepala{Bagian}\\
1 & Tentang modul ini\\\garis
2 & Masalah pembuka: Siapa yang di atas rata-rata?\\\garis
3 & Langkah 1: Deklarasi dan indeks\\\garis
4 & Langkah 2: Mengisi array dari masukan\\\garis
5 & Langkah 3: Mencari maksimum dan posisinya\\\garis
6 & Langkah 4: Dua putaran: rata-rata lalu bandingkan\\\garis
7 & Langkah 5: Di luar batas\\\garis
8 & Langkah 6: Menggeser dan membalik\\\garis
9 & Pesan error yang sering muncul\\\garis
10 & Latihan mandiri\\\garis
11 & Rangkuman\\\garis
12 & Cek pemahaman\\\garis
13 & Exit Ticket dan tugas\\\garis
14 & Bacaan lanjut dan persiapan minggu depan\\
\garisdasar
\end{tabular}}\par\vspace{6pt}
\Bagian{Pembuka}{Tentang modul ini}
Modul ini mendampingi Pertemuan 9, awal paruh kedua semester. Sampai UTS, setiap variabel menyimpan satu nilai. Minggu ini kalian memakai array untuk menyimpan banyak nilai sejenis di bawah satu nama, sehingga data bisa diproses berulang kali.

Isinya sama dengan slide di kelas, tetapi ditulis lebih pelan dan lebih lengkap, supaya kalian bisa mengulang sendiri di rumah tanpa harus menebak maksud slide.

\Sub{Setelah menyelesaikan modul ini, kalian bisa}
\par\vspace{4pt}\noindent{\fontsize{9.5pt}{13pt}\selectfont
\begin{tabular}{@{}>{\raggedright\arraybackslash}p{0.070\textwidth}>{\raggedright\arraybackslash}p{0.680\textwidth}>{\raggedright\arraybackslash}p{0.230\textwidth}@{}}
\kepala{No} & \kepala{Kemampuan} & \kepala{Dibahas di}\\
1 & Menulis program yang menyimpan data sejenis di array satu dimensi dan memprosesnya dengan perulangan: mengisi, menampilkan, menjumlahkan, mencari maksimum, dan menggeser (Sub-CPMK 9.1, CPMK0607) & Langkah 1 sampai 6\\\garis
2 & Menelusuri isi array sesudah setiap langkah pemrosesan dan menjelaskan kesalahan indeks, termasuk akses di luar batas (Sub-CPMK 9.2, CPMK0608) & kotak Tebak dulu, Latihan 3\\
\garisdasar
\end{tabular}}\par\vspace{6pt}
Karena setiap instrumen penilaian dibagi rata antara Bagian A (menulis kode) dan Bagian B (menelusuri dan menjelaskan), yang dinilai bukan hanya program kalian jalan, tetapi juga kalian bisa menjelaskan mengapa program itu jalan.

\Sub{Cara memakai modul ini}
Setiap langkah punya pola yang sama: penjelasan konsep, kadang contoh dari kehidupan sehari-hari, lalu kode yang bisa langsung kalian ketik. Sesudah kode ada kotak \textbf{Tebak dulu}. Tulis tebakan kalian di kertas sebelum membaca \textbf{Jawaban} di bawahnya. Kebiasaan menebak lalu membuktikan inilah yang membuat konsep menempel, bukan sekadar membaca.

\begin{Penting}[Ketik, jangan salin]
Ketik ulang setiap contoh di GDB Online atau editor kalian, jangan disalin-tempel. Saat mengetik, kalian akan membuat salah ketik, membaca pesan error, dan memperbaikinya. Itu bagian dari belajar.
\end{Penting}
\Sub{Yang perlu disiapkan}
Peramban dengan akses ke onlinegdb.com, atau g++ di komputer kalian sendiri. Kertas dan alat tulis untuk menebak keluaran dan membuat trace table. Folder hands-on pertemuan ini dari repo mata kuliah.

\Bagian{Pembuka}{Masalah pembuka: Siapa yang di atas rata-rata?}
Dosen ingin tahu berapa mahasiswa yang nilainya di atas rata-rata kelas.

Dengan pola sentinel minggu 5, rata-rata baru diketahui setelah semua nilai dibaca. Tetapi saat itu nilai-nilai sebelumnya sudah hilang, karena setiap \texttt{cin} menimpa variabel yang sama.

\begin{MKeluaran}[title={Tampilan yang diinginkan}]
Nilai: 70 85 60 90 75 80
Rata-rata   : 76.67
Di atas rata: 3
\end{MKeluaran}
\begin{Tebak}[Coba pikirkan]
Apa yang perlu dilakukan supaya semua nilai masih ada saat rata-rata sudah diketahui? Berapa variabel yang dibutuhkan untuk 40 mahasiswa?
\end{Tebak}
\begin{Jawaban}[Cara memecahnya]
Simpan semua: Semua nilai perlu disimpan, bukan hanya totalnya. Proses dua kali: Putaran pertama menghitung rata-rata, putaran kedua membandingkan. Terjemahkan: Satu array, dua \texttt{for} dengan indeks 0 sampai n − 1.
\end{Jawaban}
\Bagian{Langkah 1}{Deklarasi dan indeks}
Array adalah sederet variabel bertipe sama yang berbagi satu nama. Setiap elemen diakses dengan indeks di dalam kurung siku, dimulai dari 0. Array \texttt{int nilai[5]} punya elemen \texttt{nilai[0]} sampai \texttt{nilai[4]}.

Setiap elemen berperilaku seperti variabel biasa: bisa dibaca, ditugaskan, dan dipakai di ekspresi.

\begin{MKode}{array.cpp}
int nilai[5] = {70, 85, 60, 90, 75};
cout << nilai[0] << " " << nilai[4] << endl;
nilai[2] = 65;
nilai[3] += 5;
for (int i = 0; i < 5; i++) {
    cout << nilai[i] << " ";
}
cout << endl;
\end{MKode}
\begin{Tebak}[]
Sebelum menjalankan program, tulis keluarannya di kertas, karakter demi karakter.
\end{Tebak}
\begin{Jawaban}[]
\begin{MKeluaran}[]
70 75
70 85 65 95 75 
\end{MKeluaran}
Array berukuran 5 punya indeks 0 sampai 4. Ukuran harus konstanta.
\end{Jawaban}
\Bagian{Langkah 2}{Mengisi array dari masukan}
Bila banyak data baru diketahui saat program berjalan, deklarasikan array dengan ukuran maksimum yang masuk akal, lalu simpan banyak data sebenarnya di variabel \texttt{n}. Semua perulangan memakai batas \texttt{i < n}.

\begin{MKode}{baca.cpp}
int n, a[100];
cin >> n;
for (int i = 0; i < n; i++) {
    cin >> a[i];
}
int total = 0;
for (int i = 0; i < n; i++) {
    total += a[i];
}
cout << "Total " << total << endl;
\end{MKode}
\begin{MKeluaran}[title={Masukan}]
4
10 20 30 40
\end{MKeluaran}
\begin{Tebak}[]
Sebelum menjalankan program, tulis keluarannya di kertas, karakter demi karakter.
\end{Tebak}
\begin{Jawaban}[]
\begin{MKeluaran}[]
Total 100
\end{MKeluaran}
Deklarasikan array cukup besar (misalnya 100), lalu pakai hanya n elemen pertama.
\end{Jawaban}
\Bagian{Langkah 3}{Mencari maksimum dan posisinya}
Untuk mencari maksimum, anggap elemen pertama terbesar, lalu bandingkan dengan elemen berikutnya satu per satu. Menyimpan indeks lebih berguna daripada menyimpan nilai, karena indeks juga memberi tahu posisinya.

\begin{MKode}{maks.cpp}
int a[6] = {70, 85, 60, 90, 75, 80};
int iMaks = 0;
for (int i = 1; i < 6; i++) {
    if (a[i] > a[iMaks]) {
        iMaks = i;
    }
}
cout << a[iMaks] << " di indeks " << iMaks << endl;
\end{MKode}
\begin{Tebak}[]
Sebelum menjalankan program, tulis keluarannya di kertas, karakter demi karakter.
\end{Tebak}
\begin{Jawaban}[]
\begin{MKeluaran}[]
90 di indeks 3
\end{MKeluaran}
Simpan indeks, bukan nilainya. Dari indeks, nilai selalu bisa diambil kembali.
\end{Jawaban}
\Bagian{Langkah 4}{Dua putaran: rata-rata lalu bandingkan}
Inilah jawaban masalah pembuka. Putaran pertama menghitung rata-rata. Putaran kedua membandingkan setiap nilai dengan rata-rata itu. Tanpa array, putaran kedua mustahil dilakukan karena nilai-nilainya sudah hilang.

\begin{MKode}{diatas.cpp}
int a[6] = {70, 85, 60, 90, 75, 80};
int total = 0;
for (int i = 0; i < 6; i++) total += a[i];
double rata = total / 6.0;
int diAtas = 0;
for (int i = 0; i < 6; i++) {
    if (a[i] > rata) diAtas++;
}
cout << rata << " " << diAtas << endl;
\end{MKode}
\begin{Tebak}[]
Sebelum menjalankan program, tulis keluarannya di kertas, karakter demi karakter.
\end{Tebak}
\begin{Jawaban}[]
\begin{MKeluaran}[]
76.6667 3
\end{MKeluaran}

\end{Jawaban}
\Bagian{Langkah 5}{Di luar batas}
C++ tidak memeriksa apakah indeks berada di dalam array. Membaca \texttt{a[3]} dari array berukuran 3 tidak menghasilkan error kompilasi; program membaca memori di sebelah array. Hasilnya bisa acak, bisa sama di komputer kalian tetapi berbeda di komputer lain, atau program berhenti mendadak.

Kebiasaan yang mencegahnya: selalu tulis syarat perulangan sebagai \texttt{i < n}, bukan \texttt{i <= n}.

\begin{MKode}{batas.cpp}
int a[3] = {1, 2, 3};
int jumlah = 0;
for (int i = 0; i < 3; i++) {
    jumlah += a[i];
}
cout << jumlah << endl;
// for (int i = 0; i <= 3; i++) membaca a[3]:
// di luar array, hasilnya tidak bisa ditebak
\end{MKode}
\begin{Tebak}[]
Sebelum menjalankan program, tulis keluarannya di kertas, karakter demi karakter.
\end{Tebak}
\begin{Jawaban}[]
\begin{MKeluaran}[]
6
\end{MKeluaran}
C++ tidak memeriksa batas indeks. Program bisa tetap jalan dengan hasil acak, atau berhenti mendadak.
\end{Jawaban}
\Bagian{Langkah 6}{Menggeser dan membalik}
Operasi yang memindahkan elemen, seperti rotasi, paling mudah dipahami dengan menelusuri isi array setelah setiap putaran. Perhatikan arah pergeseran: menggeser ke kanan harus dimulai dari belakang. Bila dimulai dari depan, nilai di depan menimpa nilai berikutnya sebelum sempat dipindah.

\par\vspace{4pt}\noindent{\fontsize{9.5pt}{13pt}\selectfont
\begin{tabular}{@{}>{\raggedright\arraybackslash}p{0.196\textwidth}>{\raggedright\arraybackslash}p{0.414\textwidth}>{\raggedright\arraybackslash}p{0.370\textwidth}@{}}
\kepala{Langkah} & \kepala{Isi array} & \kepala{Keterangan}\\
awal & 10 20 30 40 50 & simpan = 50\\\garis
i = 4 & 10 20 30 40 40 & a[4] = a[3]\\\garis
i = 3 & 10 20 30 30 40 & a[3] = a[2]\\\garis
i = 1 & 10 10 20 30 40 & setelah i = 2 dan i = 1\\\garis
akhir & 50 10 20 30 40 & a[0] = simpan\\
\garisdasar
\end{tabular}}\par\vspace{6pt}
\Bagian{Waspada}{Pesan error yang sering muncul}
Semua pesan di tabel ini dihasilkan g++ dengan opsi \texttt{-Wall}, sama seperti yang dipakai \texttt{cek.sh}.
\par\vspace{4pt}\noindent{\fontsize{9.5pt}{13pt}\selectfont
\begin{tabular}{@{}>{\raggedright\arraybackslash}p{0.240\textwidth}>{\raggedright\arraybackslash}p{0.270\textwidth}>{\raggedright\arraybackslash}p{0.470\textwidth}@{}}
\kepala{Kesalahan} & \kepala{Contoh} & \kepala{Pesan pertama dari g++}\\
Ukuran negatif & \texttt{int a[-3];} & \texttt{error: narrowing conversion of '-3' from 'int' to 'long unsigned int'}\\\garis
Inisialisasi kelebihan elemen & \texttt{int a[2] = \{1, 2, 3\};} & \texttt{error: too many initializers for 'int [2]'}\\\garis
Indeks bukan bilangan bulat & \texttt{a[1.5]} & \texttt{error: invalid types 'int [3][double]' for array subscript}\\\garis
Menugaskan array ke array & \texttt{b = a;} & \texttt{error: invalid array assignment}\\\garis
Ukuran tidak ditulis & \texttt{int a[];} & \texttt{error: storage size of 'a' isn't known}\\
\garisdasar
\end{tabular}}\par\vspace{6pt}
Ukuran array dari variabel, misalnya \texttt{int a[n]} dengan \texttt{n} dibaca dari \texttt{cin}, diterima g++ tanpa pesan apa pun sebagai ekstensi, tetapi itu bukan C++ standar dan ditolak compiler lain. Di mata kuliah ini, ukuran array selalu konstanta.

Akses di luar batas, misalnya \texttt{a[5]} pada \texttt{int a[5]} atau \texttt{a[i]} dengan \texttt{i} = n, tidak dilaporkan saat kompilasi. Itulah sebabnya tracing indeks penting.

\Bagian{Latihan}{Latihan mandiri}
Latihan ini sama dengan hands-on di kelas. Semua file ada di \texttt{02 Hands-on/P9 - Array 1D - Hands-on/latihan/}. Kerjakan berurutan, lalu cocokkan dengan \texttt{expected/} atau jalankan \texttt{bash cek.sh}.
\par\vspace{4pt}\noindent{\fontsize{9.5pt}{13pt}\selectfont
\begin{tabular}{@{}>{\raggedright\arraybackslash}p{0.070\textwidth}>{\raggedright\arraybackslash}p{0.360\textwidth}>{\raggedright\arraybackslash}p{0.150\textwidth}>{\raggedright\arraybackslash}p{0.400\textwidth}@{}}
\kepala{No} & \kepala{File} & \kepala{Tingkat} & \kepala{Yang dilatih}\\
1 & \texttt{handson1-nilai.cpp} & Dasar & mengisi dan menampilkan, urutan terbalik\\\garis
2 & \texttt{handson2-statistik.cpp} & Menengah & rata-rata, maksimum, dua putaran\\\garis
3 & \texttt{handson3-geser.cpp} & Tantangan & telusuri isi array, perbaiki tiga kesalahan indeks\\\garis
+ & \texttt{bonus-frekuensi.cpp} & Pengayaan & array sebagai pencacah\\
\garisdasar
\end{tabular}}\par\vspace{6pt}
\Sub{Latihan 1: handson1-nilai.cpp}
Baca n nilai ke array, tampilkan dengan indeksnya, lalu dalam urutan terbalik.

\Sub{Latihan 2: handson2-statistik.cpp}
Hitung rata-rata, nilai tertinggi beserta indeksnya, dan banyak nilai di atas rata-rata.

\Sub{Latihan 3: handson3-geser.cpp}
Rotasi kanan yang salah indeks. Gambar isi array sesudah setiap putaran, perbaiki, dan jelaskan.

\Sub{Bonus: bonus-frekuensi.cpp}
Histogram frekuensi nilai per rentang dengan array pencacah.

\Sub{Latihan menelusuri}
\begin{MKode}{tebak.cpp}
int a[5] = {3, 1, 4, 1, 5};
for (int i = 1; i < 5; i++) {
    a[i] = a[i] + a[i - 1];
}
for (int i = 0; i < 5; i++) {
    cout << a[i] << " ";
}
cout << endl;
\end{MKode}
\begin{Tebak}[]
Tulis keluarannya di kertas tanpa menjalankan program. Buat trace table bila perlu.
\end{Tebak}
\begin{Jawaban}[]
\begin{MKeluaran}[]
3 4 8 9 14 
\end{MKeluaran}
Setiap elemen ditambah elemen sebelumnya yang sudah diperbarui, sehingga terbentuk jumlah kumulatif: 3, 4, 8, 9, 14.
\end{Jawaban}
\Bagian{Penutup}{Rangkuman}
\par\vspace{4pt}\noindent{\fontsize{9.5pt}{13pt}\selectfont
\begin{tabular}{@{}>{\raggedright\arraybackslash}p{0.320\textwidth}>{\raggedright\arraybackslash}p{0.660\textwidth}@{}}
\kepala{Konsep} & \kepala{Yang perlu diingat}\\
Deklarasi dan indeks & Array berukuran 5 punya indeks 0 sampai 4.\\\garis
Mengisi array dari masukan & Deklarasikan array cukup besar (misalnya 100), lalu pakai hanya n elemen pertama.\\\garis
Mencari maksimum dan posisinya & Simpan indeks, bukan nilainya.\\\garis
Dua putaran: rata-rata lalu bandingkan & Inilah jawaban masalah pembuka.\\\garis
Di luar batas & C++ tidak memeriksa batas indeks.\\\garis
Menggeser dan membalik & Operasi yang memindahkan elemen, seperti rotasi, paling mudah dipahami dengan menelusuri isi array setelah setiap putaran.\\
\garisdasar
\end{tabular}}\par\vspace{6pt}
\Bagian{Penutup}{Cek pemahaman}
Jawab tanpa menjalankan program, lalu periksa dengan menjalankannya.
\begin{enumerate}
\item Apa indeks elemen terakhir dari \texttt{double d[12]}?
\item Mengapa maksimum sebaiknya disimpan sebagai indeks?
\item Apa yang terjadi bila program membaca \texttt{a[10]} dari \texttt{int a[10]}?
\item Mengapa menghitung banyak nilai di atas rata-rata butuh array?
\item Rotasi kanan harus menggeser dari arah mana, dan mengapa?
\end{enumerate}
\begin{Jawaban}[]
\begin{enumerate}[leftmargin=16pt]
\item 11.
\item Indeks memberi nilai sekaligus posisinya.
\item Perilaku tidak terdefinisi; hasil acak atau program berhenti.
\item Rata-rata baru diketahui setelah semua nilai dibaca, jadi nilai harus disimpan untuk dibandingkan.
\item Dari belakang, agar tidak menimpa nilai yang belum dipindah.
\end{enumerate}
\end{Jawaban}
\Bagian{Penilaian}{Exit Ticket 9}
Dikerjakan di 25 menit terakhir kelas, sendiri, tanpa AI, internet, atau diskusi. Exit Ticket 9 adalah satu dari 14 Exit Ticket; rata-ratanya menjadi komponen berbobot 25\%.
\par\vspace{4pt}\noindent{\fontsize{9.5pt}{13pt}\selectfont
\begin{tabular}{@{}>{\raggedright\arraybackslash}p{0.080\textwidth}>{\raggedright\arraybackslash}p{0.600\textwidth}>{\raggedright\arraybackslash}p{0.100\textwidth}>{\raggedright\arraybackslash}p{0.200\textwidth}@{}}
\kepala{Soal} & \kepala{Isi} & \kepala{Poin} & \kepala{CPMK}\\
A1 & Tulis program yang membaca n nilai ke array dan menampilkan nilai minimum beserta indeksnya & 25 & CPMK0607\\\garis
A2 & Tulis program yang menghitung banyak elemen genap dan ganjil & 25 & CPMK0607\\\garis
B1 & Telusuri isi array setelah sebuah perulangan yang mengubah elemennya & 20 & CPMK0608\\\garis
B2 & Jelaskan mengapa sebuah perulangan menyebabkan akses di luar batas & 20 & CPMK0608\\\garis
B3 & Jelaskan mengapa sebuah masalah butuh array, bukan variabel tunggal & 10 & CPMK0608\\
\garisdasar
\end{tabular}}\par\vspace{6pt}
\Bagian{Penilaian}{Tugas 9: Sistem Nilai Mahasiswa}
Kumpulkan sebelum Pertemuan 10 lewat kanal pengumpulan yang diumumkan dosen.

\Sub{Bagian A: program (50 poin)}
Buat program dengan ketentuan berikut. Gunakan \texttt{const} untuk ukuran array.
\begin{enumerate}
\item Baca jumlah mahasiswa (maksimal 50)
\item Baca NIM dan nilai (0 sampai 100) setiap mahasiswa ke dua array sejajar, dengan validasi
\item Tampilkan semua data, nilai tertinggi dan terendah, serta rata-rata kelas
\item Tampilkan jumlah LULUS (nilai 60 ke atas), TIDAK LULUS, dan persentase kelulusan
\item Fitur pencarian: baca sebuah NIM, tampilkan nilainya atau pesan tidak ditemukan
\end{enumerate}
\begin{Contoh}[Bonus, tidak wajib]
Urutkan nilai dari tertinggi ke terendah (materi sorting di Pertemuan 12).
\end{Contoh}
\Sub{Bagian B: penjelasan (50 poin)}
Komentar maksimal 150 kata dan gambar isi kedua array setelah membaca tiga data contoh. Jelaskan cara menjaga indeks tetap di dalam batas, dan mengapa dua array harus diproses dengan indeks yang sama.

\Sub{Rubrik Tugas 9}
\par\vspace{4pt}\noindent{\fontsize{8.5pt}{11.5pt}\selectfont
\begin{tabular}{@{}>{\raggedright\arraybackslash}p{0.190\textwidth}>{\raggedright\arraybackslash}p{0.070\textwidth}>{\raggedright\arraybackslash}p{0.200\textwidth}>{\raggedright\arraybackslash}p{0.180\textwidth}>{\raggedright\arraybackslash}p{0.170\textwidth}>{\raggedright\arraybackslash}p{0.170\textwidth}@{}}
\kepala{Kriteria} & \kepala{Poin} & \kepala{Sangat baik 85–100} & \kepala{Baik 70–84} & \kepala{Cukup 55–69} & \kepala{Kurang <55}\\
A1 Program berjalan & 20 & tanpa error dan peringatan & ada peringatan & perlu satu perbaikan & tidak terkompilasi\\\garis
A2 Kelengkapan fitur & 15 & semua statistik dan pencarian & satu fitur kurang & dua kurang & tiga atau lebih kurang\\\garis
A3 Validasi dan ketepatan & 15 & validasi tepat, hasil benar untuk n = 1 dan n = 10 & satu salah & dua salah & sebagian besar salah\\\garis
B1 Isi array dan batas indeks & 30 & gambar tepat, alasan jelas & satu langkah keliru & gambar tanpa alasan & tidak ada atau disalin\\\garis
B2 Cerita error dan pengujian & 20 & pesan, sebab, perbaikan, uji & pesan tidak dikutip & hanya perbaikan & tidak ada\\
\garisdasar
\end{tabular}}\par\vspace{6pt}
Nilai kriteria = poin × skor level. Contoh: A1 di level Baik dengan skor 80 memberi 20 × 0,80 = 16 poin.

\Bagian{Penutup}{Bacaan lanjut dan persiapan minggu depan}
\begin{enumerate}
\item Deitel, P. dan Deitel, H. C++ How to Program, edisi ke-10, Bab 7.
\item learncpp.com, Bab 17 (C-style arrays).
\item Gaddis, T. Starting Out with C++, Bab 7.
\item Slide \texttt{01 Slide/P9 Array Satu Dimensi.pdf} dan hands-on di \texttt{02 Hands-on/P9 - Array 1D - Hands-on/}.
\end{enumerate}
\begin{Penting}[Minggu depan]
Pertemuan 10 membahas array dua dimensi dan string, sejalan dengan Algoritma Pertemuan 10.
\end{Penting}
\end{document}
